
Now these two squares are the ``base-change'' squares obtained from the two squares (7.10.0), so that one can apply (IV 3). One need only verify that these squares correspond to pairs @@M0@@ and @@M1@@ of schemes over @@M2@@ that are Tor-independent over @@M3@@. This is trivial for the first pair, since @@M4@@ is flat, and also immediate for the second, since the immersion @@M5@@ is regular of codimension 1 and remains so after the base change @@M6@@.

Now no longer suppose that @@M7@@ is proper, but suppose instead that it has finite Tor-dimension. I claim that the same is then true of @@M8@@ and @@M9@@. This is trivial for the first, and for the second follows from the fact (just noted) that @@M10@@ and @@M11@@ are Tor-independent over @@M12@@. One can therefore consider the square of homomorphisms
\sgaStableTarget{sga6.e10.eq.7-10-2}@@M13@@
where the horizontal arrows are defined in (IV 2.12). Note that, when @@M14@@ and @@M15@@ satisfy the conditions of 7.7, i.e., have regular fibers, @@M17@@ automatically has finite Tor-dimension (since @@M16@@ is regular), and the preceding diagram can be interpreted as a commutative diagram of rings @@M18@@.

To prove the commutativity of (7.10.2), it again suffices to prove that of the front squares of the prism analogous to @@M19@@, but with the direct-image homomorphisms @@M20@@, @@M21@@, @@M22@@ replaced by the inverse-image homomorphisms @@M23@@, @@M24@@, @@M25@@. Now the commutativity of these squares follows at 